\section{三角函数的定义}
\subsection{任意角}
初中阶段，所谓的角是两条射线共顶点形成的图形，它的大小范围是0°到360°。现在，为了方便函数处理，我们重新给角下一个定义以扩展它的范围到$(-\infty,+\infty)$：
\begin{definition}[角]
    角是由两条射线共顶点、其中一条射线固定（称为角的始边），另外一条射线（称为角的终边）逆时针旋转形成的图形。逆时针旋转经过的角度称为角的大小，如果顺时针旋转则认为角的大小是负值。图\ref{angle}展示了一个大小为$\theta$的角。
    \begin{figure}[H]
        \centering
        \begin{tikzpicture}
            \coordinate (origo) at (0,0);
            \draw[thick,->] (origo) -- ++(4,0) node[black,right] (xaxis){$x$};
            \draw[thick,gray,->] (origo) -- ++(0,4) node (mary) [black,above] {$y$};
            \draw[thick] (origo) -- ++(300:3) coordinate (bob);
            \pic [draw, ->, red, angle eccentricity=2] {angle = xaxis--origo--bob} node[left, color=red] {$\theta$};
        \end{tikzpicture}
        \caption{角}
        \label{angle}
    \end{figure}
\end{definition}
此外，为了方便计算弧长，我们重新定义角度的单位：
\begin{definition}[弧度制]
    角度制下，在半径为$r$的圆中，大小为$x$的圆心角对应的弧长为
    $$
        l=\frac{x}{360}2\pi r
    $$
    其中$360$是周角对应的大小。现在我们重新定义一个角度单位，使得角的大小为$\alpha$的时候
    $$
        l = \alpha r
    $$
    那么，周角对应的角度大小就是$2\pi$。这样的角单位制称为弧度制，单位写作$\mathrm{rad}$。
\end{definition}
\subsection{三角比}
\begin{definition}[有限角三角比]
    在直角三角形中，我们定义了一系列的三角比。具体来说有：
    \begin{enumerate}
        \item 正弦：$$\sin \theta  = \frac{\text{对边}}{\text{斜边}}$$
        \item 余弦：
              $$
                  \cos \theta  = \frac{\text{邻边}}{\text{斜边}}
              $$
        \item 正切：
              $$
                  \tan \theta  = \frac{\text{对边}}{\text{邻边}}
              $$
              以及其他的不是很常用的三角比：
        \item 正割：
              $$
                  \sec \theta = 1/\cos \theta
              $$
        \item 余割：
              $$
                  \csc \theta = 1/\sin \theta
              $$
        \item 余切：
              $$
                  \cot \theta = 1/\tan \theta
              $$
    \end{enumerate}
\end{definition}
在将角推广到旋转角后，角的大小就是任意值了。因此原有的三角比定义不再成立，需要做相应的推广。
\begin{definition}[任意角三角比]
    \begin{figure}[H]
        \centering
        \begin{tikzpicture}[scale=1.5]
            \draw (0,0) circle (2);
            \coordinate (origo) at (0,0);
            \coordinate[label = above left:$P$] (P) at (1,{sqrt(3)});
            \draw[thick,->] (origo) -- ++(3,0) node[black,right] (xaxis){$x$};
            \draw[thick,gray,->] (origo) -- ++(0,3) node (mary) [black,above] {$y$};
            \draw[thick] (origo) -- ++(60:2) coordinate (bob);
            \pic [draw, ->, red, angle eccentricity=2] {angle = xaxis--origo--bob} node[left, color=red] {$\theta$};
            \draw[blue, ->] (1,{sqrt(3)}) -- (1,0);
            \draw[cyan, ->] (1,{sqrt(3)}) -- (0,{sqrt(3)});
            \node[color=cyan] at (-0.5,{sqrt(3)}) {$\cos \theta$};
            \node[color=blue] at (1,-0.3) {$\sin \theta$};
            \draw[black,dotted] (-2,2)--(3,2);
            \draw[black,dotted] ({2*sqrt(3)/3},2)--(1,{sqrt(3)});
            \draw[brown,->] ({2*sqrt(3)/3},2)--({2*sqrt(3)/3},0);
            \node[color=brown] at (1.6,1) {$\tan \theta$};
            \node at (P)[circle,fill,inner sep=2pt]{};
        \end{tikzpicture}
        \caption{三角函数线}
    \end{figure}
    如图所示，直角坐标系内大小为$\theta$的角的终边和单位圆的交点为$P$，分别作出图中的三条\textbf{有向线段}即为三角函数线。

    他们的大小由$P(x,y)$的坐标决定：
    $$
        \sin \theta = x
    $$
    $$
        \cos \theta = y
    $$
    $$
        \tan \theta = \frac{x}{y}
    $$
    如此一般，即便角大于$\pi$也可以表示出它的三角比。
\end{definition}
\begin{definition}[三角函数]
    使用上述三角函数线的定义，我们有：

    正弦函数：
    $$
        y=\sin x \quad x\in \mathbb{R}
    $$
    余弦函数：
    $$
        y=\cos x \quad x\in \mathbb{R}
    $$
    正弦函数：
    $$
        y=\tan x=\frac{\sin x}{\cos x} \quad \cos x \ne 0
    $$
    统称为三角函数。
\end{definition}
\subsection{诱导公式}
根据任意角三角比的三角函数线定义，我们知道三角函数是存在周期性（旋转一周，所有三角函数长度方向都不变）的：
$$
    \sin (x+2\pi) = \sin x
$$
$$
    \cos (x+2\pi) = \cos x
$$
$$
    \tan (x+2\pi) = \tan x
$$
半周期（旋转半周，所有三角函数线长度不变，方向相反）则会使三角函数变为相反数：
$$
    \sin (x+\pi) = -\sin x
$$
$$
    \cos (x+\pi) = -\cos x
$$
$$
    \tan (x+\pi) = -\tan x
$$
并且根据坐标定义，不难得出他们的奇偶性：
$$
    \sin (-x) = -\sin x
$$
$$
    \cos (-x) = \cos x
$$
$$
    \tan (-x) = -\tan x
$$
最后，根据坐标定义。$\theta$和它的余角$\frac{\pi}{2}-\theta$的终边是关于$\frac{\pi}{4}$对称的，从而横纵坐标交换的。它们的三角函数存在以下关系：
$$
    \sin (\frac{\pi}{2} - x) = \cos x
$$
$$
    \cos (\frac{\pi}{2} - x) = \sin x
$$
\section{同角三角关系}
根据任意角三角比的坐标定义，显然有：
\begin{enumerate}
    \item 商数关系： $$\tan x = \frac{\sin x}{\cos x}$$
    \item 平方关系：$$\sin^2x+\cos^2x=1$$
\end{enumerate}

\section{异角三角关系}
异角三角关系的一切基于下面的定理：
\begin{theorem}[两角差的余弦]
    $$\cos(a-b) = \cos a\cos b+\sin a+\sin b$$
\end{theorem}
\begin{proof}
    $a-b$终边和单位圆的交点记作$P(\cos(a-b), \sin(a-b))$，$a$和$b$的终边和单位圆的交点分别记作$A(\cos a,\sin a)$和$B(\cos b, \sin b)$。并且记$(1,0)$点为T。那么根据图像的旋转不变性显然有
    $$AB=PT$$
    带入坐标之间的距离公式得到：
    $$
        \sqrt{(\cos a-\cos b)^2+(\sin a-\sin b)^2}=\sqrt{(\cos(a-b)-1)^2+\sin^2(a-b)}
    $$
    展开整理得到：
    $$
        2-2\cos a\cos b-2\sin a\sin b = 2-2\cos (a-b)
    $$
    也即是
    $$
        \cos(a-b) = \cos a\cos b+\sin a\sin b
    $$
\end{proof}
把$b$替换为$b+\frac{\pi}{2}$，再结合诱导公式我们可以得到两角差的正弦公式：
$$
    \sin (a-b) = \sin a\cos b-\cos a\sin b
$$
再根据三角函数的奇偶性，把$b$替换为$-b$可以得到两角和的正余弦公式：
$$
    \sin(a+b) = \sin a\cos b+\cos b\sin a
$$
$$
    \cos(a+b) = \cos a\cos b-\sin b\sin a
$$
从而根据商数关系得到：
$$
    \tan(a+b) = \frac{\sin a\cos b+\cos b\sin a}{\cos a\cos b-\sin b\sin a}
$$
分子分母同除以$\cos a\cos b$得到：
$$
    \tan(a+b) = \frac{\tan a+\tan b}{1-\tan a\tan b}
$$
根据奇偶性，把$b$替换为$-b$得到：
$$
    \tan(a-b) = \frac{\tan a-\tan b}{1+\tan a\tan b}
$$
\subsection{倍角公式}
二倍角公式较为常用，实际上就是两角和的正余弦和正切公式的简单应用：
\begin{enumerate}
    \item 正弦二倍角：
          $$
              \sin(2x) = 2\sin x\cos x
          $$
    \item 余弦二倍角：
          $$
              \cos(2x) = \cos^2x-\sin^2x
          $$
    \item 正切二倍角：
          $$
              \tan(2x) = \frac{2\tan x}{1-\tan^2x}
          $$
\end{enumerate}
二倍角公式把一次项转化成了二次项，所以也常叫升幂公式。与此相对的还有降幂公式：
\begin{enumerate}
    \item 交叉项降幂：
          $$
              \sin x\cos x = \frac{1}{2} \sin(2x)
          $$
    \item 平方降幂：
          $$
              \sin^2x = \frac{1-\cos (2x)}{2}
          $$
          $$
              \cos^2x = \frac{1+\cos (2x)}{2}
          $$
\end{enumerate}
最后，根据以上倍角公式可知，$x$的所有三角函数都可以由$t = \tan (x/2)$表示，这也称为万能代换：
\begin{enumerate}
    \item 正弦代换
          $$
              \sin x = \frac{2t}{1+t^2}
          $$
    \item 余弦代换
          $$
              \cos x = \frac{1-t^2}{1+t^2}
          $$
    \item 正切代换
          $$
              \tan x = \frac{2t}{1-t^2}
          $$
\end{enumerate}
这样一来所有的三角函数都可以表示为$\tan(x/2)$的有理函数，大大简化了问题。
\subsection{积化和差}
很容易验证：
$$
    \sin a\sin b = -\frac{1}{2}\left[\cos(a+b)-\cos(a-b)\right]
$$
$$
    \cos a\cos b = \frac{1}{2}\left[\cos(a+b)+\cos(a-b)\right]
$$
$$
    \sin a\cos b = \frac{1}{2}\left[\sin(a+b)+\sin(a-b)\right]
$$
以上三个公式是把三角函数的乘积转化为和差的公式，简称积化和差公式。
\subsection{和差化积}
很容易验证：
$$
    \sin a+\sin b = 2\sin \frac{a+b}{2}\cos \frac{a-b}{2}
$$
$$
    \sin a-\sin b = 2\cos \frac{a+b}{2}\sin \frac{a-b}{2}
$$
$$
    \cos a+ \cos b = 2\cos \frac{a+b}{2}\cos \frac{a-b}{2}
$$
$$
    \cos a- \cos b = -2\sin \frac{a+b}{2}\sin \frac{a-b}{2}
$$
以上三个公式是把三角函数的和差转化为乘积的公式，简称和差化积公式。
\section{正、余弦定理}
\begin{theorem}[正弦定理]
    在$\triangle ABC$中，$A$、$B$、$C$三个角对应的三边长分别为$a$、$b$、$c$。$R$是三角形外接圆的半径，那么
    $$
        \frac{a}{\sin A} = \frac{b}{\sin B} = \frac{c}{\sin C} = 2R
    $$
\end{theorem}
\begin{proof}
    在外接圆中利用同弦所对圆周角相等，此定理很容易证明。
\end{proof}

\begin{theorem}[余弦定理]
    在$\triangle ABC$中，$A$、$B$、$C$三个角对应的三边长分别为$a$、$b$、$c$。那么
    $$
        \cos A = \frac{b^2+c^2-a^2}{2bc}
    $$
\end{theorem}
\begin{proof}
    把三角形摆在直角坐标系中，利用距离公式立即得证。
\end{proof}

\section{$y=A\sin(\omega x+\phi)$}
这个函数的性质根据诱导公式以及我们前面提到的函数图像变换\ref{transform}可以很容易得出。

值得一提的是，这个函数有很强的物理背景，简谐运动的轨迹就是这个函数的图像。
\section{反三角函数}

由于三角函数的函数值和自变量都不是一一对应的，然而我们经常需要问“正弦为x的角是多大？”这样的问题。所以我们在一个特定的区间上规定反三角函数：

\begin{enumerate}
    \item 反正弦函数
          $$
              y = \arcsin x \quad x\in [-\frac{\pi}{2}, \frac{\pi}{2}]
          $$
    \item 反余弦函数
          $$
              y = \arccos x \quad x\in [0, \pi]
          $$
    \item 反正切函数
          $$
              y = \arctan x \quad x\in (-\frac{\pi}{2}, \frac{\pi}{2})
          $$
\end{enumerate}
\begin{problemset}
    \item 求$n$倍角$\sin(nx)$的一般公式
\end{problemset}